\chapter{Implementation: Chronological Integration of ML-DSA into the TDX/DCAP Attestation Stack}

This chapter describes the implementation work performed to integrate ML-DSA
into the TDX attestation stack used in this project. The discussion is limited
to the engineering changes introduced in the codebase and to the architectural
problems encountered during development. The cryptographic construction of
ML-DSA, the TDX attestation model, the DCAP verification pipeline, and the
meaning of quote collateral have already been introduced in the previous
chapters.

The implementation followed a chronological progression. The first steps
introduced an algorithm identity and a quote layout for ML-DSA. The following
steps moved the new algorithm through the public wrapper, the internal quote
logic, the TDQE ECALL interface, the TDQE sealed state, and the quote signing
path. The last steps moved the implementation from simulation to a real
bare-metal TDX host, added host--guest orchestration, and addressed the
verifier-side difference between a stock Intel TDQE identity and the modified
TDQE used by this work.

\section{Phase 1 -- Quote Layout and Algorithm Identity}

The first implementation step was to make ML-DSA representable in the quote
format. The existing code path was centered on ECDSA and did not contain an
attestation-key identifier, signature-data structure, or verifier allowlist
entry for ML-DSA.

The quote header version was not changed. The generated TDX quote remains a
version 4 quote. The post-quantum extension is represented through the
\texttt{att\_key\_type} field, not through the quote version field. The
repository-local quote headers were extended as follows:

\begin{itemize}
  \item \texttt{SGX\_QL\_ALG\_ECDSA\_P256 = 2};
  \item \texttt{SGX\_QL\_ALG\_MLDSA\_65 = 5};
  \item \texttt{SGX\_QL\_ALG\_MLDSA\_87 = 6}.
\end{itemize}

The implementation therefore treats the tuple

\begin{equation}
  (quote\_version, att\_key\_type)
\end{equation}

as:

\begin{equation}
  (4, 2) \quad \text{for ECDSA P-256},
\end{equation}

\begin{equation}
  (4, 5) \quad \text{for ML-DSA-65},
\end{equation}

and:

\begin{equation}
  (4, 6) \quad \text{for ML-DSA-87}.
\end{equation}

The ML-DSA size constants added to the quote headers were:

\begin{equation}
  |pk_{65}| = 1952,\quad |\sigma_{65}| = 3309,
\end{equation}

\begin{equation}
  |pk_{87}| = 2592,\quad |\sigma_{87}| = 4627.
\end{equation}

The TDQE-side private-key sizes used by the sealed blob implementation were:

\begin{equation}
  |sk_{65}| = 4032,\quad |sk_{87}| = 4896.
\end{equation}

New quote signature-data structures were added for version 4 quotes. Each
ML-DSA signature-data object contains the quote signature, the attestation
public key, and the certification-data area that follows the algorithm-specific
signature payload.

The first architectural problem appeared immediately on the verifier side. QVL
contained an attestation-key allowlist that accepted only the ECDSA type. As a
consequence, a quote with \texttt{att\_key\_type=5} or
\texttt{att\_key\_type=6} could not reach any ML-DSA-specific parser logic,
because it was rejected before the authentication-data layout was inspected.
This made the quote-format change a cross-cutting modification: the generation
headers, the wrapper headers, and the QVL constants had to evolve together.

A second problem was caused by header resolution. Some compilation units were
including SDK-provided quote headers rather than the modified repository-local
headers. The SDK headers did not contain the ML-DSA constants. The TDX wrapper
headers were therefore changed to include the repository-local quote headers
explicitly.

\section{Phase 2 -- Wrapper Context and Algorithm Propagation}

After the quote format could represent ML-DSA, the public TDX wrapper had to
carry the selected algorithm through the quote-generation path. The wrapper
previously behaved as an ECDSA-oriented interface: even when a caller supplied
a different attestation key identifier, the internal path still depended on
ECDSA defaults unless the selected algorithm was stored and propagated
explicitly.

The wrapper context was extended with an algorithm field:

\begin{center}
\texttt{m\_att\_key\_algorithm\_id}
\end{center}

The default value remains \texttt{SGX\_QL\_ALG\_ECDSA\_P256}. When the caller
creates a context with an ML-DSA attestation key id, the wrapper records
\texttt{SGX\_QL\_ALG\_MLDSA\_65} or \texttt{SGX\_QL\_ALG\_MLDSA\_87}. The
public wrapper then dispatches the following operations through
algorithm-aware methods:

\begin{itemize}
  \item \texttt{init\_quote};
  \item \texttt{get\_quote\_size};
  \item \texttt{get\_quote};
  \item \texttt{get\_keyid}.
\end{itemize}

The TDQE ECALL interface was changed accordingly. Both
\texttt{gen\_att\_key} and \texttt{gen\_quote} now receive an explicit
\texttt{algorithm\_id}. This prevented the selected algorithm from being lost
between the untrusted wrapper and the trusted TDQE implementation.

Unsupported algorithms are rejected through a dedicated TDQE error. This was
important because a silent fallback to ECDSA would have produced valid-looking
quotes with the wrong \texttt{att\_key\_type}. During testing, this distinction
became essential: some early hardware paths requested an ML-DSA UUID but still
returned an ECDSA quote header. The implementation therefore treats the quote
header as the final authority for whether the selected algorithm actually
survived the whole path.

\section{Phase 3 -- ML-DSA Backend Integration in the TDQE Build}

The next step was to integrate the ML-DSA backend into the TDQE build. The
implementation uses a TDQE-facing adapter rather than calling the imported
backend directly from all TDQE call sites. The adapter exposes one uniform
interface for each supported ML-DSA variant:

\begin{itemize}
  \item key generation;
  \item signing;
  \item signature verification.
\end{itemize}

For ML-DSA-65, the adapter test confirmed that a fixed seed produces a
deterministic key pair, that signing succeeds, that verification succeeds on
the original message, and that verification rejects a modified message or a
modified signature. The observed ML-DSA-65 adapter sizes were:

\begin{equation}
  |pk_{65}| = 1952,\quad |sk_{65}| = 4032,\quad |\sigma_{65}| = 3309.
\end{equation}

The backend integration exposed several concrete implementation problems. The
repository initially referenced ML-DSA source files that were not present at
the expected paths. To keep the parent repository self-contained, the ML-DSA
65 and 87 wrapper translation units were moved into the parent
\texttt{ae/pq} area rather than relying on modifications inside a nested
submodule path.

The SGX enclave toolchain also did not provide \texttt{cet.h}. The CET path
was therefore disabled only for the ML-DSA native object by compiling that
object with the CET macro undefined. This avoided changing the global enclave
build configuration.

The imported backend also exposed a randomized API path that required
\texttt{randombytes}. The TDQE integration, however, requires deterministic
key derivation from TDQE-controlled material. The backend configuration was
therefore changed to disable the randomized API for this integration path. A
host-side \texttt{randombytes} stub was used only in the standalone adapter
test, not as the TDQE source of attestation-key randomness.

Finally, the imported signing API expected the pure-ML-DSA domain-separation
prefix to be prepared by the caller. The adapter was updated to construct that
prefix before calling the backend internal signing and verification functions.
This change was necessary before the adapter could pass deterministic
sign/verify tests.

\section{Phase 4 -- ML-DSA Sealed Blob State}

After the backend worked in isolation, the TDQE persistent state had to be
extended. The original sealed blob path was ECDSA-specific. It contained an
ECDSA plaintext metadata structure, an ECDSA private key in the sealed payload,
and ECDSA-specific validation logic. ML-DSA required a parallel persistent
state without sharing the ECDSA blob.

The implementation introduced algorithm-specific persistent labels:

\begin{itemize}
  \item \texttt{tdqe\_data.blob} for ECDSA;
  \item \texttt{tdqe\_data\_mldsa\_65.blob} for ML-DSA-65;
  \item \texttt{tdqe\_data\_mldsa\_87.blob} for ML-DSA-87.
\end{itemize}

For each ML-DSA variant \(v \in \{65,87\}\), the TDQE constructs a sealed
container:

\begin{equation}
  B_v = \mathsf{Seal}_{TDQE}(P_v, C_v),
\end{equation}

where \(P_v\) is the authenticated plaintext metadata and \(C_v\) is the
encrypted secret payload. The plaintext metadata contains the ML-DSA public
key, the key identifier, authentication data, and certification-related fields.
The encrypted payload contains the ML-DSA private key.

The key identifier follows the existing attestation-key binding pattern:

\begin{equation}
  id_v = \mathsf{SHA256}(pk_v \parallel authentication\_data).
\end{equation}

The TDQE uses this identifier in the report data used by the certification
flow. This preserves the existing relationship between the attestation public
key, authentication data, and the certification state.

The main issue in this phase was that adding new structures was not sufficient.
The existing functions were not generic blob handlers. The following paths had
to be generalized:

\begin{itemize}
  \item blob-size selection;
  \item persistent label selection;
  \item blob sealing;
  \item blob unsealing;
  \item blob verification;
  \item key-id extraction;
  \item certification-data storage;
  \item quote signing.
\end{itemize}

The implementation therefore introduced algorithm-aware blob-size helpers and
a generic blob verification path capable of opening ECDSA, ML-DSA-65, and
ML-DSA-87 blobs while preserving their distinct plaintext and ciphertext
layouts.

\section{Phase 5 -- TDQE Quote Generation with ML-DSA}

The next phase connected the ML-DSA sealed blob to the actual quote generation
operation. The TDQE \texttt{gen\_quote} implementation was generalized so that
it can unseal either the ECDSA blob or one of the ML-DSA blobs, extract the
algorithm-specific private key, and write the matching quote signature-data
layout.

The implemented ML-DSA quote signing operation is:

\begin{equation}
  \sigma_v = \mathsf{Sign}_{\mathrm{MLDSA}\mbox{-}v}(sk_v, m_{quote}),
\end{equation}

where \(m_{quote}\) is the signed quote payload and \(sk_v\) is recovered only
inside the TDQE after unsealing \(B_v\). Immediately after signing, the TDQE
performs:

\begin{equation}
  \mathsf{Verify}_{\mathrm{MLDSA}\mbox{-}v}(pk_v, m_{quote}, \sigma_v).
\end{equation}

If the internal verification fails, quote generation fails. This mirrors the
existing ECDSA path, where the TDQE verifies the generated signature before
returning the quote.

The most invasive part of this phase was not the signature call itself, but
the removal of ECDSA-specific assumptions around the quote construction. The
original path assumed ECDSA for:

\begin{itemize}
  \item the sealed plaintext type;
  \item the sealed ciphertext type;
  \item the attestation public-key field;
  \item the signature field;
  \item the certification-data copy;
  \item the QE report and authentication-data pointers.
\end{itemize}

To avoid duplicating the whole quote-generation function, the implementation
uses algorithm-neutral pointers after the selected blob is unsealed. The
selected branch fills those pointers from the ECDSA, ML-DSA-65, or ML-DSA-87
plaintext structures. The shared quote-construction code then copies the
selected metadata into the final quote.

The certification-storage ECALL also had to be changed. It originally exposed
an ECDSA plaintext structure. This was replaced with an opaque byte buffer and
an explicit size. The TDQE infers the algorithm from the blob size, validates
that the plaintext size matches the selected algorithm, checks the corresponding
key identifier against the QE report data, and reseals the updated blob.

\section{Phase 6 -- Build System and Simulation Bring-Up}

Before the hardware machine was available, the modified TDQE, wrapper, QGS,
and verifier had to run in the local simulation environment. This phase was
necessary to validate the internal state transitions before attempting real
TDX hardware execution.

The build system had to be adjusted in several locations:

\begin{itemize}
  \item the TDQE build was extended with the ML-DSA adapter and native backend
        objects;
  \item the quote verification build was extended with the ML-DSA helper
        sources;
  \item the QVL CMake build was aligned with the Makefile build so that unit
        tests saw the ML-DSA verifier code;
  \item the PCE wrapper, QGS, and TDX quote wrapper were made sensitive to
        \texttt{SGX\_MODE};
  \item versioned enclave symlinks were created where the runtime expected
        \texttt{.so.1} names.
\end{itemize}

The main problems in this phase were not isolated compiler errors, but
inconsistent runtime assumptions between components. The PCE wrapper linked
against hardware SGX runtime libraries even when the rest of the stack was
using simulation libraries. QGS and the TDX quote wrapper also needed
simulation-aware runtime linking. Without this alignment, the ML-DSA bootstrap
failed before reaching the trusted quote logic.

The TDQE loader also expected the signed enclave at
\texttt{libsgx\_tdqe.signed.so.1}, while the build produced
\texttt{libsgx\_tdqe.signed.so}. This was fixed by creating the expected
versioned symlink. In addition, QGS was changed to accept the TDQE signed
enclave path explicitly through \texttt{QGS\_TDQE\_PATH}; deriving the TDQE
path from the wrapper library directory was incorrect in the repository-local
layout.

Simulation exposed one further limitation: \texttt{sgx\_verify\_report2}
caused a \texttt{SIGILL} in the simulated path. The implementation introduced
explicit \texttt{SE\_SIM} bypasses for the affected checks so that the local
simulation path could continue to exercise ML-DSA key generation, sealing, and
quote assembly. These bypasses are marked in the logs and are not used as
evidence of hardware-backed attestation.

\section{Phase 7 -- QGS Selection and Context Management}

Once the TDQE and wrapper could represent ML-DSA, QGS had to select the
requested attestation key correctly. The initial QGS behavior still depended
on the default ECDSA context. During a \texttt{GET\_QUOTE\_REQ}, this caused
the request to carry an ML-DSA key id while the initialized context remained
ECDSA.

The implementation embedded the local default attestation key identifiers used
by the test stack:

\begin{itemize}
  \item ECDSA: \texttt{...4f3f};
  \item ML-DSA-65: \texttt{...4f40};
  \item ML-DSA-87: \texttt{...4f41}.
\end{itemize}

These values are selection identifiers. They are not attestation private keys.
QGS now uses the requested identifier to initialize the selected algorithm
context before quote generation. This prevents ML-DSA requests from being
routed through the default ECDSA context.

This phase also exposed a memory-management bug. Switching from the default
ECDSA context to an ML-DSA context could free the old context twice. The QGS
cleanup path was adjusted so that context replacement does not double-free the
default context.

The final architectural outcome is that QGS remains an untrusted orchestrator:
it selects the requested algorithm and routes the request, but it does not own
or use the ML-DSA private key. The private key remains sealed and is only
unsealed inside the TDQE.

\section{Phase 8 -- Verifier and QVL Support}

After quote generation began producing ML-DSA-shaped quotes, the verifier path
had to be extended. The first verifier probe confirmed that generation was no
longer the only blocker: QVL still rejected ML-DSA before signature verification
because the parser and allowlist were ECDSA-only.

The QVL implementation was extended in four areas:

\begin{itemize}
  \item attestation-key constants and allowlists;
  \item ML-DSA v4 authentication-data parsing;
  \item generic accessors for the quote signature and attestation public key;
  \item ML-DSA signature verification in the quote verifier branch.
\end{itemize}

The verifier-side ML-DSA check verifies the quote signature over the signed
quote payload using the attestation public key carried in the quote signature
area. The local verifier fallback also checks the binding between the
attestation public key, QE authentication data, and the QE report:

\begin{equation}
  \mathsf{SHA256}(attest\_pub\_key \parallel qe\_auth\_data)
  =
  QEReport.report\_data.
\end{equation}

This fallback was added only to classify and validate repository-local
experiments when standard collateral verification could not complete. It does
not replace DCAP collateral verification.

Once the parser and ML-DSA signature branch were added, the next failure moved
to certification data. The verifier could parse and validate the ML-DSA quote
structure, but \texttt{tee\_qv\_get\_collateral} rejected local-only
certification data. The observed local-only type was \texttt{3}, whereas the
standard collateral path expects a PCK-chain-compatible certification layout.

This result was important because it separated three concerns:

\begin{itemize}
  \item quote parsing was working;
  \item ML-DSA quote signature verification was working;
  \item standard collateral compatibility was still not solved.
\end{itemize}

\section{Phase 9 -- PCCS, QCNL, and Collateral Diagnostics}

The next phase instrumented the collateral path. The implementation added
diagnostics around the QPL and QCNL calls used by quote generation and
verification. The goal was to distinguish malformed platform inputs from
missing cache entries, network failures, TLS issues, and unsupported
certification-data layouts.

The ML-DSA simulation runner was extended with a local PCCS bootstrap path.
It can generate local TLS material, start the PCCS service in a repository
controlled configuration, and point the QCNL test configuration at that local
service. Both \texttt{pccs\_url} and \texttt{collateral\_service} are set
explicitly because they are used by different parts of the collateral path.

When the local PCCS cache was empty, the ML-DSA quote generation path could
fall back to local-only certification data only when explicitly enabled by:

\begin{center}
\texttt{TDX\_MLDSA\_LOCAL\_PCCS\_EMPTY\_FALLBACK=1}
\end{center}

This fallback is intentionally scoped to the repository-local test path. It is
not enabled as a general hardware behavior.

\section{Phase 10 -- Migration to the Bare-Metal TDX Host}

The implementation was then moved to a new bare-metal TDX-capable machine.
This phase was required because the initial development machine did not expose
the device needed by direct TDX guest attestation. The new host had to be
brought from a generic Ubuntu 24.04 installation to a working TDX host with
attestation services.

The firmware configuration had to enable SGX, TME, TME-MT, TDX, the SEAM
loader, and a non-zero TDX/TME key split. After the Canonical TDX host setup
and reboot, the host showed TDX initialization in the kernel log and
\texttt{/sys/module/kvm\_intel/parameters/tdx} returned \texttt{Y}.

The host attestation stack required \texttt{qgsd} and \texttt{pccs}. PCCS was
configured for local connections and lazy cache fill. A valid Intel PCS
subscription key was required because the configured PCS endpoint was the
production endpoint:

\begin{center}
\texttt{https://api.trustedservices.intel.com/sgx/certification/v4/}
\end{center}

Before the key and platform cache were corrected, PCCS returned \texttt{401}
errors from Intel PCS, and \texttt{qgsd} logged that no certificate data was
available for the platform. The host then used \texttt{PCKIDRetrievalTool} to
populate PCCS for the platform. After this step, PCCS accepted platform
requests and \texttt{qgsd} returned successful quote initialization and quote
generation results for the stock ECDSA path.

This phase also clarified that the Canonical registration check and the
\texttt{PCKIDRetrievalTool} workflow are not the same path. The registration
script may report the MPA-oriented status, while the platform can still be
populated for PCCS through \texttt{PCKIDRetrievalTool}.

\section{Phase 11 -- Host--Guest Test Orchestration}

Hardware quote generation had to run inside a TDX guest, while QGS, PCCS, and
platform collateral remain host-side services. A host-side orchestration script
was therefore added:

\begin{center}
\texttt{tdx\_tests/direct/run\_host\_tdx\_guest\_repo\_tests.sh}
\end{center}

The script performs host preflight checks, creates or reuses a TDX guest image,
boots the guest, installs guest dependencies, synchronizes the repository into
the guest, cleans stale build artifacts, runs the selected ML-DSA and ECDSA
flows, and prints host \texttt{qgsd}/PCCS diagnostics on failure.

Several realization problems appeared in this phase. The repository copy step
initially attempted to archive local PCCS TLS and cache files, including
root-owned private-key material, which caused permission failures. Those
generated artifacts were excluded from the copy. The copied build tree also
contained stale dependency files with absolute paths from the host, causing
guest-side \texttt{make} failures. These stale files had to be removed before
guest builds.

Another problem was a broken absolute symlink in the repo-local SGX SDK:
\texttt{libsgx\_urts.so.2} pointed to a path from the previous machine. The
runner repairs this symlink in the guest. Finally, the hardware guest path
could not safely inherit all repo-local SGX library paths. In hardware mode,
using incompatible repo-local runtime libraries conflicted with the system PSW
and caused quote initialization failures.

\section{Phase 12 -- Hardware ML-DSA Quote Generation}

Once the host and guest were operational, the ML-DSA runner was changed from a
simulation-first structure to a hardware-first structure. When
\texttt{/dev/tdx\_guest} is present, the runner executes the direct TDX
hardware path.

The verifier probe and the direct capability probe both select the requested
ML-DSA key id through the direct TDX attestation interface and check the final
quote header. The test fails if the selected key id is all-zero, if the
selected key id does not match the requested ML-DSA identifier, or if the quote
header does not contain the expected ML-DSA attestation key type.

The stock host \texttt{qgsd} path worked for ECDSA but did not preserve the
ML-DSA selection end to end in this repository-local setup. In earlier runs,
the ML-DSA UUID was requested, but the returned quote still had
\texttt{att\_key\_type=2}. Forcing a guest-local repo-local QGS failed for a
different reason: the guest-local path still required SGX QE/uRTS/PSW support
inside the guest and failed during enclave loading.

The working hardware design was therefore:

\begin{itemize}
  \item keep stock host \texttt{qgsd} for ECDSA;
  \item run the patched repository-local QGS on the host for ML-DSA;
  \item expose the host repository-local QGS to the guest through vsock;
  \item run the quote request inside the TDX guest through
        \texttt{/dev/tdx\_guest}.
\end{itemize}

The observed hardware results were:

\begin{itemize}
  \item ML-DSA-65 selected key id \texttt{...4f40} and produced
        \texttt{att\_key\_type=5};
  \item ML-DSA-87 selected key id \texttt{...4f41} and produced
        \texttt{att\_key\_type=6};
  \item the local-certification ML-DSA-65 quote size observed by the direct
        probe was \texttt{7102};
  \item the ML-DSA-87 quote size observed by the direct probe was
        \texttt{9060};
  \item when the PCK-chain certification path was used for ML-DSA-65, the
        verifier probe observed quote size \texttt{10376}.
\end{itemize}

All successful hardware ML-DSA runs were reported with \texttt{SGX\_MODE=HW}
and a real TDX guest device.

\section{Phase 13 -- DCAP Verification and TDQE Identity}

The final phase addressed the distinction between quote generation,
cryptographic quote verification, DCAP collateral retrieval, and TDQE identity
verification.

The verifier probe first executes the standard DCAP API sequence:

\begin{enumerate}
  \item generate the ML-DSA quote;
  \item call \texttt{tee\_qv\_get\_collateral};
  \item call \texttt{tdx\_qv\_verify\_quote}.
\end{enumerate}

After the ML-DSA quote carried PCK-chain-compatible certification data,
\texttt{tee\_qv\_get\_collateral} could complete. The next failure was:

\begin{equation}
  \texttt{tdx\_qv\_verify\_quote} \rightarrow \texttt{0xe026}.
\end{equation}

The error \texttt{0xe026} is \texttt{QEIDENTITY\_MISMATCH}. This was not a
quote header problem and not an ML-DSA signature problem. The quote was being
generated by the modified repository-local TDQE, while the stock DCAP
collateral path supplied the stock Intel TDQE identity. The verifier correctly
rejected the quote because the embedded QE report did not match the stock Intel
TDQE identity.

Two verification modes were therefore separated.

First, a stock Intel identity mode can be requested. In that mode, the verifier
requires the standard Intel TDQE identity and fails when the modified TDQE does
not match it.

Second, a private TDQE identity mode was implemented for the repository-local
ML-DSA TDQE. In this mode, the verifier derives a local TDQE identity from the
quote's embedded QE report. The generated identity contains the report-derived
fields needed by the identity check and is explicitly marked with:

\begin{center}
\texttt{"localtdqe":true}
\end{center}

The QVL identity verifier accepts this identity only when the environment
explicitly enables the override:

\begin{center}
\texttt{TDX\_MLDSA\_ALLOW\_LOCAL\_TDQE\_IDENTITY=1}
\end{center}

and the identity JSON contains the local marker. This keeps the exception
narrow. It does not globally disable TDQE identity verification.

The implementation was then refined into a private PCS-compatible path. The
local proxy forwards Intel PCS-compatible collateral requests upstream, but
serves the TDX QE identity endpoint from the locally generated TDQE identity
file. The runner writes a QCNL configuration whose
\texttt{collateral\_service} points to this local proxy. This allows the
verification flow to keep the DCAP API surface while using the private TDQE
identity authority required by the modified TDQE.

\section{Final Implementation Result}

At the end of the implementation work, the system supports the following
properties.

\begin{itemize}
  \item The quote format contains explicit ML-DSA-65 and ML-DSA-87 attestation
        key types.
  \item The public wrapper can select ML-DSA contexts and propagate the
        selected algorithm into the TDQE ECALL surface.
  \item The TDQE can generate ML-DSA attestation keys, seal them into
        algorithm-specific blobs, unseal them later, and sign quotes with the
        private key inside the TDQE.
  \item The QGS path can select ML-DSA-65 or ML-DSA-87 contexts instead of
        falling back to the default ECDSA context.
  \item QVL can parse ML-DSA v4 quotes and verify the ML-DSA quote signature.
  \item A real TDX guest can generate hardware-backed ML-DSA quotes through
        the host repository-local QGS.
  \item The verifier can use the DCAP API path with a private TDQE identity
        authority for the modified repository-local TDQE.
\end{itemize}

The final limitation is also precise. The implementation does not make the
modified ML-DSA TDQE accepted by the stock Intel TDQE identity. The stock Intel
identity path is expected to reject the modified TDQE with
\texttt{QEIDENTITY\_MISMATCH}. The implemented verification path is therefore
best described as a DCAP API based verification path with a private TDQE
identity authority, not as unmodified stock Intel DCAP identity verification.
